\chapter{Admin Portal (web)}

Administrators sign in to the same web portal as teachers, but see one
extra tab: \textbf{Administration}. An admin can do everything a teacher
can, plus manage the school's structure and accounts.

% ----------------------------------------------------------------------
\section{Classes \& sections (fully modular)}

The school's structure is never fixed. At any time you can:

\begin{itemize}[itemsep=2pt]
  \item \textbf{Add a class}: type the name (``Grade 8'', ``BSCS-1'')
        and click \textbf{Add class}.
  \item \textbf{Add a section}: choose the class, type the section name
        (``A'', ``B'', ``Blue'') and click \textbf{Add section}.
  \item \textbf{Delete a class or section}: click \textbf{delete} next
        to it. Deletion is ``soft'' -- old attendance and reports are
        kept for the record, the class simply disappears from all
        lists.
\end{itemize}

% ----------------------------------------------------------------------
\section{Rooms \& RFID readers}

Register every physical location and the reader hardware on its door:

\begin{enumerate}[itemsep=2pt]
  \item \textbf{Add the room}: name it and pick its type -- Gate,
        Classroom, Laboratory, Library, DiscussionRoom, Auditorium or
        Other.
  \item \textbf{Add its reader}: type the reader code printed on the
        device (e.g.\ \texttt{GATE-01}), pick the room, and choose the
        number of antennas: \textbf{3 for gates, 2 for rooms}.
\end{enumerate}

From that moment the system understands the door: antenna order
$1\rightarrow2(\rightarrow3)$ means entry, the reverse means exit.

\begin{figure}[H]
\centering
\includegraphics[width=\textwidth]{admin-structure}
\caption{The Administration tab: classes \& sections on the left, rooms
\& RFID readers on the right, account creation below.}
\end{figure}

% ----------------------------------------------------------------------
\section{Creating accounts}

Nobody registers themselves -- the admin creates every account. In
\textbf{Create accounts}:

\subsection{Create a parent}
Choose account type \textbf{Parent}, fill full name, username, password
and phone, click \textbf{Create account}. Hand the username and password
to the parent.

\subsection{Create a teacher}
Same steps with account type \textbf{Teacher}.

\subsection{Create a student (and link everything together)}
Choose account type \textbf{Student}. Extra fields appear:

\begin{itemize}[itemsep=2pt]
  \item \textbf{Reg.\ no} -- the student's registration number;
  \item \textbf{RFID tag EPC} -- the code of the RFID chip inside the
        student's ID card (your card printer or a handheld reader shows
        this code). This is what makes automatic attendance work;
  \item \textbf{Section} -- which class and section the student joins;
  \item \textbf{Parent} -- pick the parent account created earlier.
        This link is what routes gate alerts and summaries to the right
        parent's phone.
\end{itemize}

\notebox{If a student's ID card is lost, issue a new card and simply
update the RFID tag EPC on the student -- history stays intact.}

% ----------------------------------------------------------------------
\section{Semesters and the timetable}

Create a semester (name, start and end date, ``current'' flag), then
build each section's timetable in the \textbf{Timetable} tab: day,
start/end time, subject, teacher, room -- entry by entry. Parents and
students see changes in their apps immediately.

\begin{figure}[H]
\centering
\includegraphics[width=\textwidth]{admin-timetable}
\caption{Building a section timetable as admin.}
\end{figure}

% ----------------------------------------------------------------------
\section{Daily \& weekly summary times}

By default parents receive the daily summary at \textbf{18:00} and the
weekly summary on \textbf{Friday}. Your IT person can change both in
the server settings (\texttt{Summaries:DailyTime},
\texttt{Summaries:WeeklyDay}).

% ----------------------------------------------------------------------
\section{Checklist: setting up a new school}

\begin{enumerate}[itemsep=2pt]
  \item Create the semester.
  \item Add classes and their sections.
  \item Add rooms, then the RFID reader of each room (3 antennas at
        gates, 2 elsewhere).
  \item Create teacher accounts.
  \item Create parent accounts.
  \item Create student accounts -- linking each student to a section, a
        parent, and the RFID EPC of their ID card.
  \item Build the timetable for every section.
  \item Ask parents and students to install the app and sign in.
\end{enumerate}

\tipbox{The default administrator login shipped with the system is
\texttt{admin} / \texttt{Admin@123} -- change this password on day
one.}
